\section{Логистична регресия. Linear Discriminant \\ analysis.}

	\subsection{Кратка информация}
	Мотивация - логистичната регресия се използва за Credit Score модели.\\
	Моделът има вида: \\
	$\mathbb{E}(Y|X_1, \dots, X_p) = Prob(Y=1|X_1, \dots, X_p) = \frac{e^{\beta_0 + \beta_1 X_1 + \dots + \beta_p X_p}}{1+ e^{\beta_0 + \beta_1 X_1 + \dots + \beta_p X_p}} $ \\
	$\mathbb{E}(Y|X_1, \dots, X_p) = Prob(Y=0|X_1, \dots, X_p) = \frac{e^{\beta_0 + \beta_1 X_1 + \dots + \beta_p X_p}}{1+ e^{\beta_0 + \beta_1 X_1 + \dots + \beta_p X_p}} $
	
	Идея за други класификационни модели - Linear Discriminant Analysis, Quadratic Discriminant Analysis, 
	Linear Discriminant analysis \\
	Ако класовете са с многомерно гаусово разпределение, или
	$$f_k(x) = \frac{1}{(2\pi)^p/2 \left | \Sigma_k  \right |} e^{-\frac{1}{2} (x-\mu_k)^T \Sigma_k^{-1}(x-\mu_k)} $$
	в специалния случай, където ковариационните матрици са равни, $\Sigma_k = \Sigma \forall k$, имаме 
	Линеен дискриминантен анализ.
	
	
	\subsection{Упражнения}
	Conceptual 4.8 section, ex 1, 2, 12 \\

	\begin{exercise}[Sec 4.8, Conceptual ex. 1]
			Using a little bit of algebra, prove that (4.2) is equivalent to (4.3). In
		other words, the logistic function representation and logit represen-
		tation for the logistic regression model are equivalent.
	\end{exercise}

	\begin{exercise}[Sec 4.8, Conceptual ex. 2]
		А
	\end{exercise}


	\begin{exercise}[Sec 4.8, Conceptual ex. 12]
	Suppose that you wish to classify an observation $x \in R $ into apples
	and oranges. You fit a logistic regression model and find that
	
	Pr(Y = orange|X = x) =
	%exp(ˆβ0 + ˆβ1x)
	%1 + exp(ˆβ0 + ˆβ1x)
	.
	Your friend fits a logistic regression model to the same data using the
	softmax formulation in (4.13), and finds that
	$$Pr(Y= orange| X=x) = \frac{exp(\hat{a}_{orange} + \hat{a}_{orangex}) }{den}$$
	Pr(Y = orange|X = x) =
	%exp(ˆαorange0 + ˆαorange1x)
	%exp(ˆαorange0 + ˆαorange1x) + exp(ˆαapple0 + ˆαapple1x).
	(a) What is the log odds of orange versus apple in your model?
	(b) What is the log odds of orange versus apple in your friend’s
	model?
	(c) Suppose that in your model, $\beta_0$ and $\beta_1$. What are
	the coefficient estimates in your friend’s model? Be as specific
	as possible.
	(d) Now suppose that you and your friend fit the same two models
	on a different data set. This time, your friend gets the coefficient
	%estimates ˆαorange0 = 1.2, ˆαorange1 = −2, ˆαorange0 = 3, ˆαorange1 =
	0.6. What are the coefficient estimates in your model?
	(e) Finally, suppose you apply both models from (d) to a data set
	with 2,000 test observations. What fraction of the time do you
	expect the predicted class labels from your model to agree with
	those from your friend’s model? Explain your answer. \\
	\end{exercise}


	
	
	
	\newpage